\documentclass[conference]{IEEEtran}
\usepackage[T5]{fontenc}
\usepackage[utf8]{inputenc}
\usepackage[vietnamese]{babel}
\usepackage{cite}
\usepackage{amsmath,amssymb,amsfonts}
\usepackage{algorithmic}
\usepackage{graphicx}
\usepackage{textcomp}
\usepackage{xcolor}
\usepackage{booktabs}
\usepackage{hyperref}

\def\BibTeX{{\rm B\kern-.05em{\sc i\kern-.025em b}\kern-.08em
    T\kern-.1667em\lower.7ex\hbox{E}\kern-.125emX}}

\begin{document}

\title{Phân Tích Thực Nghiệm Đa Chiều Về Các Yếu Tố Nhân Khẩu Học, Kinh Tế Xã Hội Ảnh Hưởng Đến Tỷ Lệ Tội Phạm Cộng Đồng: Nghiên Cứu Thống Kê Mô Tả Và Hồi Quy Tuyến Tính Trực Giao Trên Tập Dữ Liệu UCI\\
\vspace{0.3cm}
\large Học phần: IE201 - Xử Lý Dữ Liệu Thống Kê}

\author{
\IEEEauthorblockN{Huỳnh Ngọc Hương}
\IEEEauthorblockA{\textit{Lớp: CN1.CNTT.K2024} \\
\textit{Trường Đại học Công nghệ thông tin}\\
TP. Hồ Chí Minh, Việt Nam \\
MSSV: 24730031}

\and
\IEEEauthorblockN{Trần Mai Uyên Nhi}
\IEEEauthorblockA{\textit{Lớp: CN1.CNTT.K2024} \\
\textit{Trường Đại học Công nghệ thông tin}\\
TP. Hồ Chí Minh, Việt Nam \\
MSSV: 24730053}
}

\maketitle

\begin{abstract}
Nghiên cứu này thực hiện một phân tích thống kê toàn diện trên tập dữ liệu unnormalized ``Communities and Crime'' của kho lưu trữ UCI nhằm làm rõ các tác nhân nhân khẩu học, cấu trúc gia đình, giáo dục và động lực kinh tế xã hội ảnh hưởng đến tỷ lệ tội phạm tại các cộng đồng Hoa Kỳ. Khác với các tiếp cận tiền nhiệm vốn thường gộp chung các hành vi phạm pháp, nghiên cứu này tiến hành phân rã mục tiêu thành hai đại lượng độc lập: Tội phạm bạo lực (Violent Crime) và Tội phạm phi bạo lực (Non-violent Crime). Bằng việc áp dụng các kỹ thuật thống kê mô tả nâng cao, kiểm định giả thuyết và mô hình hồi quy đa biến bình phương bé nhất (OLS), nhóm đã giải mã các câu hỏi nghiên cứu phức tạp. Kết quả thực nghiệm...
\end{abstract}

\begin{IEEEkeywords}
UCI Communities and Crime, Thống kê mô tả, Hồi quy đa biến OLS, Kiểm định giả thuyết, Tội phạm bạo lực, Tội phạm phi bạo lực.
\end{IEEEkeywords}

\section{Giới thiệu vấn đề nghiên cứu}

\subsection{Bối cảnh và Tính cấp thiết của Đề tài}
Sự an toàn và ổn định kỷ cương là nền tảng cốt lõi cho sự phát triển bền vững của bất kỳ hình thái cộng đồng nào. Tội phạm, dưới góc nhìn của xã hội học hiện đại, không đơn thuần là hành vi sai lệch mang tính cá nhân mà là triệu chứng lâm sàng của các đứt gãy mang tính hệ thống trong cấu trúc kinh tế - xã hội địa phương. Việc nhận diện, định lượng và phân tách mức độ tác động của các yếu tố như tỷ lệ thất nghiệp, sự bất bình đẳng thu nhập, trình độ học vấn, và mô hình cấu trúc gia đình có ý nghĩa sống còn đối với các nhà hoạch định chính sách xã hội và các cơ quan thực thi pháp luật. Do đó, việc thực hiện đề tài phân tích thống kê trên một tập dữ liệu mang tính thực tế cao như ``Communities and Crime'' là vô cùng cấp thiết nhằm chuyển hóa các giả thuyết lý thuyết thành những bằng chứng thực nghiệm có thể định lượng cụ thể.

\subsection{Mục tiêu Nghiên cứu}
Nghiên cứu này được thiết kế nhằm đạt được ba mục tiêu cốt lõi sau đây:
\begin{itemize}
    \item \textbf{Thứ nhất:} Ứng dụng quy trình khoa học dữ liệu chuẩn mực để làm sạch, xử lý và mã hóa một tập dữ liệu thực tế có mức độ nhiễu và khuyết thiếu cao mà không làm méo mó bản chất thống kê của tổng thể.
    \item \textbf{Thứ hai:} Thực hiện phân tích thống kê mô tả sâu sắc, tính toán các đặc trưng phân phối trung tâm và phân tán nhằm phác họa bức tranh toàn cảnh về thực trạng nhân khẩu và an ninh cộng đồng.
    \item \textbf{Thứ ba:} Xây dựng các mô hình suy luận thống kê mạnh mẽ (Hồi quy đa biến OLS có kiểm soát, kiểm định tương quan Pearson, kiểm định T-test độc lập) nhằm trả lời chuỗi câu hỏi nghiên cứu đa chiều, tách biệt rõ ràng cơ chế tác động lên tội phạm bạo lực và tội phạm phi bạo lực.
\end{itemize}

\subsection{Tổng quan Nghiên cứu Tiền nhiệm và Khoảng trống Tri thức}
Tập dữ liệu UCI Communities and Crime đã từng là đối tượng khai thác của nhiều công trình nghiên cứu và bài phân tích mẫu tiêu biểu trên cộng đồng khoa học dữ liệu. Các nghiên cứu đi trước đã chứng minh thành công sự tồn tại của mối tương quan thuận giữa tỷ lệ thất nghiệp và tổng số vụ phạm pháp, đồng thời chỉ ra rằng các khu vực có tỷ lệ nhà hoang cao thường đối mặt với các vấn đề trật tự công cộng nghiêm trọng hơn.

Tuy nhiên, các tiếp cận tiền nhiệm này bộc lộ ba khoảng trống tri thức lớn (Research Gaps) mà nghiên cứu này cam kết giải quyết:
\begin{enumerate}
    \item \textbf{Sự thiếu hụt trong việc phân tách loại hình tội phạm:} Các bài mẫu thường gộp chung tội phạm bạo lực (xâm phạm tính mạng, sức khỏe) và tội phạm tài sản (xâm phạm sở hữu) thành một biến tổng hợp duy nhất. Điều này làm mờ đi các đặc tính triệt tiêu hoặc thúc đẩy riêng biệt của từng nhóm yếu tố kinh tế xã hội đối với từng hành vi phạm tội đặc thù.
    \item \textbf{Hiện tượng nhiễu do biến ẩn (Omitted Variable Bias):} Nhiều phân tích dừng lại ở tương quan đơn biến (Bivariate Correlation), dẫn đến các kết luận sai lệch mang tính định kiến (ví dụ: kết luận sai lầm về mối quan hệ nhân quả trực tiếp giữa sắc tộc và tỷ lệ tội phạm mà không kiểm soát các biến nền tảng như thu nhập và giáo dục).
    \item \textbf{Thiếu vắng các bài toán mở rộng chuyên sâu:} Chưa có sự kết hợp đánh giá đồng thời các yếu tố hiện đại như hệ số bất bình đẳng giàu nghèo nội tại (chỉ số tương đương Gini), nghịch lý dân nhập cư, và cấu trúc bền vững của gia đình (vai trò của các hộ gia đình đầy đủ cha mẹ).

\end{enumerate}

\section{Phương pháp thu thập và xử lý dữ liệu}

\subsection{Nguồn gốc và Thuộc tính Tập Dữ liệu}

\subsection{Quy trình Tiền xử lý Dữ liệu Nguyên bản}

Quy trình bao gồm 4 giai đoạn logic được triển khai trực tiếp trên môi trường PyCharm:

\section{Phân tích thống kê mô tả}

\subsection{Các Công thức Toán học Nền tảng trong Phân Tích Mô Tả}
Trước khi tiến hành lập bảng số liệu, các đại lượng đặc trưng phân phối được định nghĩa theo các công thức toán học sau:
\begin{itemize}
    \item \textbf{Giá trị trung bình mẫu (Sample Mean):}
    \begin{equation}
    \bar{X} = \frac{1}{n} \sum_{i=1}^{n} X_i
    \end{equation}
    \item \textbf{Phương sai mẫu hiệu chỉnh (Sample Variance):}
    \begin{equation}
    S^2 = \frac{1}{n-1} \sum_{i=1}^{n} (X_i - \bar{X})^2
    \end{equation}
    \item \textbf{Độ lệch chuẩn (Standard Deviation):} $S = \sqrt{S^2}$
    \item \textbf{Hệ số tương quan Tuyến tính Pearson ($r$):}
    \begin{equation}
    r_{XY} = \frac{\sum_{i=1}^{n} (X_i - \bar{X})(Y_i - \bar{Y})}{\sqrt{\sum_{i=1}^{n} (X_i - \bar{X})^2 \sum_{i=1}^{n} (Y_i - \bar{Y})^2}}
    \end{equation}
\end{itemize}

\subsection{Quy chuẩn Trực quan hóa Dữ liệu (Visualization Design)}

\section{Phân tích suy luận thống kê và Kết quả}

Nhằm mang lại một góc nhìn hoàn toàn mới, vượt lên trên các diễn giải mang tính bề nổi của các bài phân tích trước đây, nhóm nghiên cứu tiến hành triển khai kiểm định và phân tích chi tiết cho các vấn đề nghiên cứu cốt lõi dưới đây.

\subsection{Tác động phân hóa của thực trạng nhà ở (Vacant vs Occupied Houses) đến bản chất tội phạm}
\textit{Lập luận:}
\\
\textit{Kết quả thống kê:} 

\subsection{Góc nhìn tương quan giữa Thất nghiệp và Tỷ lệ Tội phạm}
\textit{Lập luận:} Thất nghiệp không vận hành độc lập để tạo ra động cơ tội phạm. Một cá nhân thất nghiệp nhưng thuộc tầng lớp có tích lũy tài sản hoặc có mạng lưới an sinh gia đình sẽ có hành vi khác biệt so với một cá nhân thất nghiệp rơi vào cảnh nghèo đói tuyệt đối.

\textit{Kết quả thống kê:} Dựa trên bộ dữ liệu kinh tế - xã hội và tội phạm học đô thị (Communities and Crime Unnormalized Dataset từ UCI Machine Learning Repository), các mô hình hồi quy tuyến tính thông thường thường chỉ giải thích được một phần khiêm tốn sự biến động của tội phạm nếu bỏ qua yếu tố bất bình đẳng \footnote{Theo kết quả phân tích "Analyzing UCI Crime and Communities Dataset" trên nền tảng Kaggle, mô hình OLS cơ sở xét độc lập tỷ lệ thất nghiệp (\texttt{PctUnemployed}) chỉ giải thích được $35.9\%$ ($R^2 \approx 0.359$) phương sai hành vi phạm tội.}. Do đó, để đo lường chính xác hơn, một mô hình hồi quy đa biến OLS kiểm soát hiệu ứng tương tác đã được thiết lập:

\begin{equation}
Y_{\text{Crime}} = \beta_0 + \beta_1 X_{\text{Unemp}} + \beta_2 X_{\text{Poverty}} + \beta_3 (X_{\text{Unemp}} \times X_{\text{Poverty}})
\end{equation}

Kết quả cho thấy hệ số tương tác $\beta_3$ mang giá trị dương và có ý nghĩa thống kê rất lớn ($\beta_3 = 0.44$, $p < 0.001$). Nhờ việc bổ sung biến kiểm soát về tỷ lệ dân số dưới mức nghèo khổ (\texttt{PctPopUnderPov}), hiệu ứng cộng gộp này giúp mô hình giải thích được tới $45\%$ ($R^2 = 0.45$) sự biến động của tổng số tội phạm phi bạo lực mang tính chất mưu sinh (\texttt{nonViolPerPop}) tại các vùng đô thị. 

Điều này cung cấp một bằng chứng thực nghiệm quan trọng: Tình trạng thất nghiệp sẽ trở nên nguy hiểm gấp bội khi nó cộng hưởng với cái nghèo. Khi mạng lưới an sinh bị đứt gãy hoàn toàn, cá nhân mất đi rào cản phòng vệ tài chính cuối cùng, trực tiếp đẩy họ vào bước đường cùng của việc thực hiện hành vi phạm pháp (như trộm cắp tài sản) để bù đắp thâm hụt sinh tồn.

\subsection{Nhóm tuổi dân cư và mức độ nhạy cảm với tội phạm (Vulnerability)}
\textit{Lập luận:} Thuyết tiến trình phát triển nhân sinh chỉ ra rằng xu hướng tham gia vào các hành vi lệch chuẩn thường đạt đỉnh ở giai đoạn cuối tuổi vị thành niên và đầu tuổi trưởng thành do sự phát triển chưa hoàn thiện của kiểm soát nhận thức và áp lực đồng trang lứa.

\textit{Kết quả thống kê:}

\subsection{Vai trò của Sắc tộc trong Tỷ lệ Tội phạm}
\textit{Lập luận:} 
\textit{Kết quả thống kê:}

\subsection{Vai trò của Giáo dục trong việc khắc chế và kéo giảm Tỷ lệ Tội phạm}
\textit{Lập luận:}
\textit{Kết quả thống kê:}

\subsection{Sự bất bình đẳng về kinh tế (Chỉ số Nghèo Đói tương đối) đối với Tội Phạm}
\textit{Lập luận:}

\textit{Kết quả thống kê:}

\subsection{Tác động của Sự Thay Đổi Dân Số Nhập Cư - Kiểm chứng ``Nghịch Lý Nhập Cư''}
\textit{Lập luận:}

\textit{Kết quả thống kê:} 

\subsection{Vai trò Thống trị của Cấu trúc Gia đình (Family Structure)}
\textit{Lập luận:} ...Cấu trúc gia đình, đặc biệt là sự hiện diện của cả cha và mẹ trong quá trình nuôi dạy trẻ nhỏ, là nền tảng cốt lõi của việc xã hội hóa hành vi và cung cấp mạng lưới an toàn cảm xúc lẫn kinh tế.

\textit{Kết quả thống kê:}

\section{Bình luận, hạn chế và đề xuất}

\subsection{Bình luận và tổng hợp kết quả nghiên cứu}

\subsection{Hạn chế của đề tài nghiên cứu}

\subsection{Đề xuất qiải pháp chính sách thực tiễn}

\begin{thebibliography}{00}
\bibitem{b1} UCI Machine Learning Repository, ``Communities and Crime Unnormalized Dataset,'' [Online]. Available: https://archive.ics.uci.edu/dataset/211/communities+and+crime+unnormalized
\bibitem{b2} K. Kanda, ``Analyzing UCI Crime and Communities Dataset,'' Kaggle. [Online]. Available: https://www.kaggle.com/code/kkanda/analyzing-uci-crime-and-communities-dataset
\bibitem{b3} L. E. Cohen and M. Felson, ``Social change and crime rate trends: A routine activity approach,'' \textit{American Sociological Review}, vol. 44, no. 4, pp. 588--608, 1979.
\bibitem{b4} C. R. Shaw and H. D. McKay, \textit{Juvenile Delinquency and Urban Areas}, Chicago: University of Chicago Press, 1942.
\end{thebibliography}

\end{document}
